% =====================================================================
%  chunks/ch03j.tex —— 原书 p.125–133（PDF 132–140）
%  第 3 章 §9 配边与上配边理论作为广义同调与广义上同调（上）
% =====================================================================
% [原文 p.125]
\section[配边与上配边理论。 Adams--Novikov 谱序列。 形式群]{配边与上配边理论作为广义同调与广义上同调。配边中的上同调运算。Adams--Novikov 谱序列。形式群。循环群与圆周在流形上的作用}

\takeaway{配边论把任意光滑流形及其到空间 $X$ 的映射当作「闭链」，因此是几何上最自然的广义同调论；本节给出配边群、相对配边群与对偶的上配边群的定义，并用 Thom 空间把上配边群实现为稳定同伦类群。}

从几何观点看，\emph{配边论是最自然的同调论}。我们先从经典情形出发：其中的闭链取为任意光滑流形。在\emph{无定向配边论} $\Omega_*^O(\cdot)=\mathfrak{N}_*(\cdot)$ 中，空间 $X$ 的一个\emph{奇异 $n$-闭链}是一对 $(M^n,f)$，其中 $M^n$ 是闭流形，$f$ 是由该流形到 $X$ 的映射：
\[
(M^n,f);\qquad f:M^n\to X.
\]
% [Tutor Note: 原书在 §9 中让 bordism 与 cobordism 两词并立：bordism 指以光滑流形为「闭链」的同调侧理论 $\Omega_*^G(\cdot)$，cobordism 指与之对偶的上同调侧理论 $\Omega_G^*(\cdot)$ 以及单点的配边群；本译本在明确对举处译作「配边／上配边」，其余语境仍作「配边」，这不是原书笔误。]
两个奇异 $n$-闭链 $(M_1^n,f_1)$、$(M_2^n,f_2)$ 称为\emph{等价的}（或\emph{配边的}）：$(M_1^n,f_1)\sim(M_2^n,f_2)$，如果存在流形 $L^{n+1}$，其边界为不交并 $\partial L^{n+1}=M_1^n\cup M_2^n$，并存在映射 $g:L^{n+1}\to X$，使得
\[
g|_{M_1^n}=f_1,\qquad g|_{M_2^n}=f_2.
\]
这里假定所涉及的流形都嵌入 $\R^N$ 中，$N\gg n$。我们现定义\emph{$n$ 维配边群}为
\[
\Omega_n^O(X)=\{n\text{-闭链}\}/\sim,
\]
其中（阿贝尔）群结构按 $n$-闭链的不交并定义：
\[
[(M_1^n,f_1)]+[(M_2^n,f_2)]=[(M_1^n\cup M_2^n,\ f_1\cup f_2)].
\]
对空间对 $Y\subset X$，\emph{相对配边群} $\Omega_n^O(X,Y)$ 类似地定义：改用流形的映射 $(M^n,f)$，此时 $M^n$ 可以带边界 $\partial M^n\ne\varnothing$，边界 $\partial M$ 映入 $Y$，并适当地定义等价关系 $\sim$。

一般情形下，对任意李群的稳定序列 $G=(G_n)$（$G_n\subset O_n$），有相应的配边论 $\Omega_*^G(X,Y)$。基本例子是
\[
G=O,\ SO,\ U,\ SU,\ Sp,\ e,
\]
其中 $G_n=O_n$、$G_n=SO_n$、$G_{2n}=U_n$、$G_{2n}=SU_n$、$G_{4n}=Sp_n$、$G_n=1$。除此之外，考虑表示
\[
\rho:G\to O,\qquad \rho=(\rho_n),\qquad \rho_n:G_n\to O_n,
\]
还可产生更多例子，例如 $G=Spin=(G_n)=(Spin_n)$、$\rho_n:Spin_n\to SO_n$（二重覆叠）。

若流形 $M$ 的稳定法丛带有 $G$-结构，则称 $M$ 为 \emph{$G$-流形}。配边论 $\Omega_*^G(\cdot)$ 如上定义，只不过这里假定所涉及的流形都嵌入 $\R^N$（$N\gg n$），且其法丛带有 $G$-结构。配边群 $\Omega_n^G(X,Y)$ 具备广义同调论的全部基本性质：同伦不变性、函子性、对 CW 复形对 $Y\subset X$ 的切除公理 $\Omega_n^G(X,Y)=\Omega_n^G(X/Y,*)$，以及空间对 $(X,Y)$ 的正合列的存在性：
% [原文 p.126]
\[
\cdots\to\Omega_n^G(Y)\to\Omega_n^G(X)\to\Omega_n^G(X,Y)\to\Omega_{n-1}^G(Y)\to\cdots.
\]
单点空间 $*$ 的配边群 $\Omega_n^G(*)$ 与球面的配边群重合；这可由悬挂同构
\[
\Sigma:\Omega_n^G(X,Y)\xrightarrow{\ \cong\ }\Omega_{n+1}^G(\Sigma X,\Sigma Y)
\]
推出，据此有
\[
\Omega_n^G(*)\cong\Omega_n^G(*\cup *',*')\xrightarrow{\ \Sigma\ }\Omega_{n+k}^G(S^k,*'').
\]
单点空间的配边群称为\emph{上配边群}（cobordism groups）。它们的直和 $\Omega_*^G=\bigoplus_n\Omega_n^G(*)$，连同由流形直积诱导的乘法运算，构成\emph{上配边环}。第 4 章 §3 将指出配边群 $\Omega_n^G(*)$ 与万有向量 $G_n$-丛（底空间为 $BG_n$）的「Thom 空间」$MG=(MG_n)$ 之间的关系。「Thom 定理」给出
\[
\Omega_m^G(*)\cong\pi_m^s(MG)=\pi_{m+k}(MG_k),\qquad k>m+1.
\]
（注意 Thom 空间 $MG_n$ 是 $(n-1)$-连通的。）特别地，当 $G=e$ 时，群 $\Omega_n^e$ 与球面的稳定同伦群重合，后者已在上文 §7 末尾讨论过。关于其他群 $G$ 的上配边环 $\Omega_*^G$ 的结果将在第 4 章 §3 给出。

与配边论 $\Omega_*^G(\cdot)$ 对偶的是\emph{上配边论}（cobordism theory）$\Omega_G^*(\cdot)$，相应的\emph{上配边群}定义为
\[
\Omega_G^m(X,*)=\Omega_{N-m-1}^G(S^N\setminus X,*'),\qquad N\to\infty,
\]
其中有限 CW 复形 $X$ 嵌入维数充分大的球面中。

按 Thom 给出的方案（见第 4 章 §3）进行，可得到如下自然同构：
\begin{equation}
\Omega_G^{-m}(X,*)\cong\lim_{n\to\infty}[\Sigma^{n+m}X,MG_n], \tag{9.1}
\end{equation}
% [Tutor Note: 式 (9.1) 中的 $\lim_{n\to\infty}$ 指随 $n$ 增大的正向（归纳）极限，不是数列极限；上配边群用负指标 $-m$ 记次，是为了让它与同调的次数约定对齐。]
其中 $[\,,\,]$ 表示映射的同伦类集合。直和 $\Omega_G^*(X,Y)=\bigoplus_m\Omega_G^m(X,Y)$ 可赋予乘法结构，使之成为斜对称环。该乘法定义如下。对上述每个 $G$，群 $G_n$ 对向量丛的直和运算「封闭」，意思是这一运算给出相应的分类空间之间的映射以及其上向量丛之间的映射：
% [原文 p.127]
\begin{equation}
\begin{gathered}
BG_{n_1}\times BG_{n_2}\to BG_{n_1+n_2},\\
BG\times BG\to BG.
\end{gathered}\tag{9.2}
\end{equation}
由此得到的万有向量丛的映射
\[
\Phi_{n_1,n_2}:\eta_{n_1}^G\times\eta_{n_2}^G\to\eta_{n_1+n_2}^G
\]
在相应的 Thom 空间之间诱导「乘法」：
\[
M\Phi_{n_1,n_2}:M(G_{n_1})\wedge M(G_{n_2})\to M(G_{n_1+n_2}),
\]
其中 $V\wedge W=V\times W/V\vee W$ 是有基点空间 $V$、$W$ 的「张量积」或 \emph{smash 积}。把映射 $M\Phi_{n_1,n_2}$ 作用于映射之积 $f$、$g$：
\[
f:\Sigma^{m_1+n_1}X\to M(G_{n_1}),\qquad g:\Sigma^{m_2+n_2}X\to M(G_{n_2}),
\]
它们经由 (9.1) 表示相应上配边群的元素，于是得到映射
\[
M\Phi_{n_1,n_2}(f\times g):\Sigma^{m_1+m_2+n_1+n_2}X\xrightarrow{\ f\times g\ }M(G_{n_1})\wedge M(G_{n_2})\xrightarrow{\ M\Phi_{n_1,n_2}\ }M(G_{n_1+n_2}),
\]
从而在相应上配边群的直和 $\Omega_G^*(X)$ 上给出所需的双线性、结合乘法：
\begin{equation}
\begin{gathered}
\Omega_G^{m_1}(X)\cdot\Omega_G^{m_1}(X)\subset\Omega_G^{m_1+m_2}(X),\\
z_1\cdot z_2=(-1)^{m_1m_2}z_2\cdot z_1.
\end{gathered}\tag{9.3}
\end{equation}
% [Tutor Note: 原文 p.127 此处第二个因子印作 $\Omega_G^{m_1}(X)$；由下一行的反交换式 $z_1z_2=(-1)^{m_1m_2}z_2z_1$ 及乘积的次数 $m_1+m_2$ 可知，第二个因子应为 $\Omega_G^{m_2}(X)$，故原书系笔误，译文照印。]

带法 $G$-结构的闭流形 $M^n$ 有一个自然的\emph{基本类} $[M^n]\in\Omega_n^G(M^n)$，由恒同映射定义：
\[
1:M^n\to M^n;\qquad [M^n]=[(M^n,1)].
\]
在上配边类与配边类之间可以定义一种\emph{卡积运算}，类似于普通同调与上同调中定义的卡积（见第 2 章 §3）：
\[
\Omega_G^m(X)\cap\Omega_k^G(X)\to\Omega_k^{G-m}(X),
\]
具有性质
\[
f^*(a)\cap b=a\cap f_*(b)
\]
其中 $f$ 是 CW 复形的映射。在 $X=M^n$ 的情形，映射
\[
a\mapsto a\cap[M^n]
\]
% [原文 p.128]
定义了 \emph{Poincaré--Atiyah 对偶同构}
\[
D:\Omega_G^k(M^n)\xrightarrow{\ \cong\ }\Omega_{n-k}^G(M^n).
\]
此同构诱导出配边类的\emph{相交}运算，它与上面定义的上配边类乘积对偶：
\[
a\circ b=D^{-1}(Da\cdot Db).
\]
与任何广义上同调论一样，上配边论中存在一个（环的）谱序列，即 \emph{Atiyah--Hirzebruch 谱序列}：
\begin{equation}
\begin{gathered}
(E_r^{p,q},d_r),\qquad d_m:E_m^{p,q}\to E_m^{p+m,q-m+1},\qquad d_m^2=0,\\
E_2^{p,q}=H^p(X;\Omega_G^q(*)),
\end{gathered}\tag{9.4}
\end{equation}
其相伴环为 $\sum_{p,q}E_\infty^{p,q}=G\Omega_G^*(X)$。
% [Tutor Note: 只有谱序列退化时才有 $G\Omega_G^*(X)\cong\Omega_G^*(X)$；一般情形下 $G\Omega_G^*(X)$ 只是滤过的相伴分次环，例 3 特意指出即使 $X$ 的整上同调无挠也未必同构。]

\para{例}
\begin{example*}{1}
当 $G=O$ 时，由 Thom 1950 年代中期的定理，群 $\Omega_G^n(X)$ 同构于 $X$ 的若干不同维数的普通 $\Z/2$-上同调群的直和。环 $\Omega_G^*(X)$ 的结构由一个标准同构给出：
\[
\Omega_G^*(X)\cong H^*(X;\Omega_G^*(*)),
\]
其中 $\Omega_G^*(*)$ 是多项式代数
\[
\Omega_G^*(*)=\Z/2[v_2,v_4,v_5,v_6,v_8,\ldots],
\]
其生成元 $v_j$ 的维数为 $j=2,4,5,\ldots$，$j\ne 2^k-1$。Atiyah--Hirzebruch 谱序列退化（$d_m=0$，$m\ge 2$，故 $E_\infty^{p,q}=E_2^{p,q}$），相伴环 $G\Omega_G^*(X)=\sum_{p,q}E_\infty^{p,q}$ 同构于 $\Omega_G^*(X)$。
\end{example*}
\begin{example*}{2}
当 $G=SO$ 时，环 $\Omega_{SO}^*(X)\otimes\Z_{(2)}$ 的结构也化归为系数取自 $\Z_{(2)}$ 的普通上同调环的结构，这里 $\Z_{(2)}$ 是分母为奇数的有理数环（即 $\Z$ 在素数 2 处的局部化）；这来自 Novikov、Rohlin 与 Wall 在 1960 年前后的结果。另一方面，对任何素数 $p\ne 2$，理论 $\Omega_{SO}^*(\cdot)\otimes\Z_{(p)}$ 都不化归为普通上同调。此时环 $\Omega_{SO}^*\otimes\Z_{(p)}$ 是多项式环 $\Z_{(p)}[x_1,x_2,\ldots]$，其生成元位于形如 $4k$ 的每个负维数上（Milnor、Novikov，1960 年前后）。
\end{example*}
\begin{example*}{3}
$G=U$ 的情形最有趣，它为应用上配边论的代数技巧提供了基础。环 $\Omega_U^*$ 是多项式环，在每个偶的负维数上恰有一个生成元（Milnor、Novikov；见第 4 章 §3）：
\[
\Omega_U^*\cong\Z[u_1,u_2,\ldots],\qquad u_j\in\Omega_U^{-2j}=\Omega_{2j}^U.
\]
% [原文 p.129]
Atiyah--Hirzebruch 谱序列只对整上同调群无挠的复形 $X$ 退化；然而即使对这样的 $X$，相伴环 $G\Omega_U^*(X)$ 也未必同构于 $\Omega_U^*(X)$。
\end{example*}
\begin{example*}{4}
环 $\Omega_{SU}^*\otimes\Z_{(p)}$（$p>2$）的结构，以及 $\Omega_{SU}^*$ 的部分 2-主挠，由 Novikov 在 1960 年代初描述。$\Omega_{SU}^*\otimes\Z_{(2)}$ 的结构随后由 Conner 与 Floyd 完全阐明（1960 年代中期）。
\end{example*}
\begin{example}{5}
辛上配边环 $\Omega_*^{Sp}$ 的情形远比上述各上配边环复杂。已知环 $\Omega_*^{Sp}$ 的全部挠都是 2-主的（Novikov，1960），而 Ray 元 $\phi_j\in\Omega_{2j-5}^{Sp}$（2 阶不可分解元）在 $\Omega_*^{Sp}$ 的一切计算中起着格外重要的作用。$\Omega_*^{Sp}$ 的加法结构已借 Adams 谱序列算到维数 120（Kochman，1980 年代中期），同一维数范围内的环结构现在也基本上由 Vershinin 借助 Adams--Novikov 谱序列（1990）确定。Vershinin 还证明：带奇点的辛流形（Ray 元）的上配边环实际上化归为上配边环 $\Omega_*^U$。Botvinnik 据此用带奇点的上配边描述了辛上配边的 Adams--Novikov 谱序列；进一步的计算导致对任意 $k$ 都探测到 $\Omega_*^{Sp}$ 中 $2^k$ 阶的非平凡元（Botvinnik、Kochman，1992）。
\end{example}

Conner 与 Floyd 的工作使人们注意到：复上配边论 $\Omega_U^*(\cdot)$ 中存在 \emph{Chern 示性类}；类似地，无定向上配边论 $\Omega_O^*(\cdot)$ 中存在 Stiefel--Whitney 类。这些取值于上配边群中的示性类，是普通上同调中 Chern 类与 Stiefel--Whitney 类的直接类比。

对 $n$ 维流形 $M^n$ 及其上向量丛 $\eta$ 而言，这意味着：与示性类 $c_j(\eta)$、$w_i(\eta)$ 对偶的 $M$ 中闭链，总可以实现为流形的像。关于示性闭链的这一性质，最早的迹象由 Gamkrelidze 给出（1950 年代初）：他证明 $\CP^n$ 中代数簇的示性闭链可由代数子簇实现。

正如第 4 章 §3 将要指出的（并见上文 §6），$\Z/2$-上同调中的 Stiefel--Whitney 示性类与 Steenrod 平方之间存在密切联系，这一联系由 Thom 的公式给出；事实上，全部 Steenrod 运算 $a\in\cA_2$ 都可表示在乘积 $\RP_1^\infty\times\cdots\times\RP_n^\infty$ 的上同调中：
\[
a\mapsto a(u),\qquad u=u_1\cdots u_n\in H^n(\RP_1^\infty\times\cdots\times\RP_n^\infty).
\]
借助这一表示，可以为上配边论 $\Omega_O^*(\cdot)$ 与 $\Omega_U^*(\cdot)$ 定义 Steenrod 运算的适当类比。

下文我们只限于复上配边论 $\Omega_U^*(\cdot)$。（理论 $\Omega_O^*(\cdot)$ 类似，只是会发生相当程度的平凡化。）在任意 CW 复形 $X$ 的上配边群 $\Omega_U^2(X)$ 中，有一个由「几何上配边」类组成的子集 $W(X)$，其定义如下。
% [原文 p.130]
由于 Thom 空间 $M(U_2)$ 可以等同于 $\CP^\infty=K(\Z,2)$（见第 4 章 §3），每个上同调类 $z\in H^2(X;\Z)$ 都对应唯一一个映射的同伦类
\[
f_z:X\to M(U_2)=\CP^\infty,
\]
使得 $f_z^*(u)=z$，其中 $u\in H^2(\CP^\infty;\Z)$ 是生成元；于是映射 $f_z$ 决定一个上配边类（几何上配边类）$[z]\in\Omega_U^2(X)$。全部几何上配边类的集合 $W(X)$ 与 $H^2(X;\Z)$ 一一对应；然而这一对应不是加性的。

对 $U$-流形 $M^{2n}$，可以得到与 $W(M^{2n})$ 对偶的一族\emph{除子}
\[
\{D[z]\}\subset\Omega_{2n-2}^U(M^{2n}),
\]
它们在 $\Omega_{2n-2}^U(M^{2n})$ 中同样不是加性的。

对 $X$ 上每个复向量丛 $\eta$，我们如下定义它的第一 Chern 示性类 $\sigma_1(\eta)\in\Omega_U^2(X)$：取为如上与普通上同调中第一 Chern 类 $c_1(\eta)$ 相对应的几何上配边类 $[c_1(\eta)]$。其余的 Chern 示性类 $\sigma_j(\eta)\in\Omega_U^{2j}(X)$（$1<j\le\dim_\C\eta$）随即由一般函子性质与 Whitney 公式确定：
\begin{equation}
\sigma(\eta_1\oplus\eta_2)=\sigma(\eta_1)\sigma(\eta_2),\qquad \sigma(\eta)=1+\sigma_1(\eta)+\sigma_2(\eta)+\cdots. \tag{9.5}
\end{equation}
$\sigma_j(\eta)$ 在普通上同调群中的像仍是普通的 Chern 示性类。Chern 示性类与复上配边论 $\Omega_U^*(\cdot)$ 中的上同调运算密切相关；后者由下列条件确定：

1. 对每个整数有限序列 $\omega=(\omega_1,\ldots,\omega_k)$，存在一个稳定的加性运算
\begin{equation}
s_\omega:\Omega_U^j(X)\to\Omega_U^{j+2|\omega|}(X),\qquad \deg s_\omega=2|\omega|,\qquad |\omega|=\sum_{i=1}^k\omega_i. \tag{9.6}
\end{equation}

2. 下列乘积公式成立：
\[
s_\omega(ab)=\sum_{(\omega',\omega'')=\omega}s_{\omega'}(a)s_{\omega''}(b).
\]

3. 运算 $s_0$ 是恒同算子：$s_0=1$；若 $a\in\Omega_U^{2j}(X)$ 是几何上配边类，则
\[
s_\omega a=\begin{cases}a^{\omega_1+1} & \text{当 }\omega=(\omega_1),\\ 0 & \text{当 }\omega=(\omega_1,\ldots,\omega_k),\ k>1.\end{cases}
\]

4. 若某个全次数 $<2n$ 的乘积 $s_{\omega'}\circ s_{\omega''}$ 在元素
\[
u_1\cdots u_n\in\Omega_U^{2n}(\CP_1^\infty\times\cdots\times\CP_n^\infty)
\]
上为零（其中 $u_j\in\Omega_U^2(\CP_j^\infty)$ 是标准生成元），则该乘积是零运算：$s_{\omega'}\circ s_{\omega''}=0$。
% [原文 p.131]
由这一条件可推出复合公式
\[
s_{\omega'}\circ s_{\omega''}(u)=\sum_\omega\lambda_\omega(\omega',\omega'')s_\omega,
\]
其中 $\lambda_\omega(\omega',\omega'')$ 是整数。由此还可推出：元素 $s_\omega(u)$（$u=u_1\cdots u_n$）具有关于 $\omega_1,\ldots,\omega_n$「对称化」的单项式形式：
\[
s_\omega(u)=\Bigl(\sum u_{i_1}^{\omega_1}\cdots u_{i_n}^{\omega_n}\Bigr)u_1\cdots u_n.
\]

运算 $s_\omega$ 称为 \emph{Landweber--Novikov 运算}。它们构成 $\Z$ 上 Landweber--Novikov（LN）代数 $S$ 的一组基；$S$ 是带显然（对角）余乘的 Hopf 代数，其定义方式与 Milnor 定义经典 Steenrod 代数所用的方案相同：对角由运算作用在两个元素之积上的「Leibniz 公式」确定（见上文）。Hopf 代数在带乘法的模上的这种作用，因此称为 \emph{Milnor 模}。

不难证明：复上配边论中全部稳定运算构成的代数 $A^U$，恰好由线性组合
\[
\sum_j\lambda_j s_{\omega(j)}
\]
组成，其中 $\lambda_j\in\Lambda=\Omega_U^*$ 具有负的维数（次数）：
\[
\lambda_j\in\Omega_U^{-2q_j}=\Lambda^{-2q_j},\qquad \deg\lambda_j=-2q_j,
\]
而元素 $s_\omega$ 具有正的次数：
\[
\deg s_\omega=2|\omega|=\sum_{j=1}^k\omega_j.
\]

代数 $A^U$ 有分次
\[
A^U=\sum_{-\infty}^{\infty}A_{2n}^U,
\]
其中 $A_{2n}^U$ 由次数为 $2n$ 的运算组成，即由线性组合
\[
\sum_j\lambda_j s_{\omega(j)},\qquad \deg\lambda_j+\deg s_{\omega(j)}=2n
\]
组成。这里允许无限形式级数，只要当 $j\to\infty$ 时 $\deg s_{\omega(j)}\to\infty$。

% [原文 p.132]
运算 $s_\omega$ 在上配边环 $\Omega_U^*=\Lambda$ 上的作用，可以用具体的流形显式描述。$\Lambda$ 的每个元素都可用流形 $\CP^1,\ldots,\CP^n,\ldots$ 的上配边类的有理系数多项式表示。切丛 $\tau(\CP^n)$ 分解如下（见第 4 章 §1）：
\[
\tau(\CP^n)+\varepsilon_\C^1=(n+1)\eta,
\]
其中 $\eta$ 是 $\CP^n$ 上的标准线丛，对此有
\[
\sigma_1(\eta)=u\in\Omega_U^2(\CP^n),
\]
\[
D(u)=[\CP^{n-1}]\in\Omega_{2n-2}^U(\CP^n).
\]
由此可知，稳定法丛 $\nu(\CP^n)$ 可以等同于 $-(n+1)\eta$。对元素 $-\eta\in K^0(\CP^n)$，上配边意义下的 Chern 示性类如下：
\[
\sigma_j(-\eta)=(-1)^ju^j,\qquad Du^j=[\CP^{n-j}].
\]
我们现如下定义示性类 $\sigma_\omega$，$\omega=(\omega_1,\ldots,\omega_k)$。考虑对称多项式
\[
Q^\omega(u_1,\ldots,u_k)=\sum_\alpha u_1^{\omega_{\alpha(1)}}\cdots u_k^{\omega_{\alpha(k)}},
\]
其中求和遍历 $k$ 个字母的对称群的全部元素 $\alpha$。多项式 $Q(u_1,\ldots,u_k)$ 可唯一地写成生成元 $\sigma_j$ 的多项式 $P^\omega(\sigma_1,\sigma_2,\ldots)$，其中 $\sigma_j$ 是 $u_1,\ldots,u_k$ 的第 $j$ 个初等对称多项式。令 $\sigma_\omega$（$\omega=(\omega_1,\ldots,\omega_k)$）表示示性类
\[
\sigma_\omega=P^\omega(\sigma_1,\sigma_2,\ldots).
\]
% [Tutor Note: 注意区分 $\sigma_\omega$ 与 $s_\omega$：前者是由对称多项式 $Q^\omega$ 定义的示性类，取值于上配边群；后者是对偶的 Landweber--Novikov 运算，作用于配边环。(9.7) 把 $s_\omega$ 在 $\CP^n$ 上的作用化归为整数 $\lambda(\omega)$，这是下面「完整描述」的关键。]
现在把类 $\sigma_\omega$ 作用于丛 $-(n+1)\eta$（它与法丛 $\nu(\CP^n)$ 稳定等价），其 Chern 类 $\bar\sigma(-(n+1)\eta)$ 如下：
\[
\bar\sigma(-(n+1)\eta)=1+\bar\sigma_1+\bar\sigma_2+\cdots=\frac{1}{(1+u)^{n+1}}.
\]
一个简单计算表明：类 $\sigma_\omega(-(n+1)\eta)\in\Omega_U^{2|\omega|}(\CP^n)$ 的对偶，是配边类 $[\CP^{n-|\omega|}]$ 乘以某个整数 $\lambda(\omega)$。这就如下确定了上配边运算 $s_\omega$：
\begin{equation}
s_\omega(\CP^n)=\lambda(\omega)[\CP^{n-|\omega|}],\qquad |\omega|=\sum_{j=1}^k\omega_j. \tag{9.7}
\end{equation}
利用下列性质，即可由 (9.7) 得到上配边运算在上配边环 $\Omega_U^*=\Lambda$ 上作用的完整描述：
% [原文 p.133]
\begin{equation}
\begin{gathered}
s_\omega(\mu)=\sum_{(\omega',\omega'')=\omega}s_{\omega'}(\mu)\cdot s_{\omega''}(\mu),\qquad \mu\in\Lambda,\\
s_\omega\circ\mu=\sum_{(\omega',\omega'')=\omega}s_{\omega'}(\mu)\circ s_{\omega''},\qquad \mu\in\Lambda,
\end{gathered}\tag{9.8}
\end{equation}
连同 $s_0=1$（对空序列 $\omega$）。（这一点由 Novikov 在 1960 年代中期建立。）用复上配边论 $\Omega_U^*(\cdot)$ 的运算代数 $A^U$ 在其系数环 $\Omega_U^*=\Lambda$ 上的作用来表示该代数，这种表示是忠实的，即若 $a\in A^U$ 在 $\Lambda$ 上平凡作用，则 $a=0$。这一性质体现了复上配边论与普通上同调之间的一个根本区别。

值得注意的是，代数 $A^U$ 可以看成两个代数的乘积关于形式级数拓扑的完备化：
\[
A^U=(\Lambda\otimes S)^T.
\]
这一乘积作为向量空间同构于这两个代数的张量积，但乘法不是通常的乘法。它们构成非交换的子代数；上文给出的、关于它们换位子的 Novikov 公式（见 (9.8)）依赖于 Landweber--Novikov 代数 $S$ 在如上所述的复配边环上的作用。这里重要的是：这一作用确定了 Hopf 代数 $S$ 上的一个「Milnor 模」，即它关于余乘满足 Leibniz 乘积公式。代数 $A^U$ 实现为 Hopf 代数 $S$ 上（左）Milnor 模 $\Lambda$ 上的算子代数，而 $\Lambda$ 又实现为作用在「属于环 $\Lambda$」的函数上的、带常系数的微分算子代数。这是 Novikov（1990 年代初）在 Hopf 代数及其上 Milnor 模理论背景下给出的一类相当一般的「算子代数」构造的特例。

1978 年 Buchstaber 与 Shokurov 注意到：与 Landweber--Novikov 代数 $S$ 自然对偶的 Hopf 代数包含「单点空间的复上配边」环 $\Lambda$：
\[
\Lambda\subset X^*,\qquad \Lambda\otimes\Q=S^*\otimes\Q.
\]
% [Tutor Note: 原文 p.133 此处印作 $\Lambda\subset X^*$。本段讨论的是 $S$ 的自然对偶 Hopf 代数 $S^*$，且 $\Lambda\otimes\Q=S^*\otimes\Q$ 表明左边也应是 $S^*$；$X^*$ 在全书记号中并无定义，故原书系笔误，应读作 $S^*$，译文照印。]
（通过 Chern 数——Milnor 与 Novikov 的早期结果给出了 $\Lambda$ 的整结构——Buchstaber 在 1970 年代后期从这一新观点分析了环 $\Lambda$；其结构原来极不平凡。）因此我们可以把这一构造看作 Hopf 代数 $S$ 及其对偶 $S^*$ 的「算子双」这一非常一般的构造（Novikov 1992 意义下）的特例。任意代数 $S$ 在其对偶空间 $S^*$ 上有自然的作用 $R_a^*$、$L_b^*$，它们分别是右乘与左乘的伴随：
\[
R_b(a)=L_a(b)=ab,\qquad (R_y^*(z),x)=(z,R_y(x))=(z,xy).
\]

% CHUNK ch03j DONE
